\chapter{Ricerca}\label{chap:dolor}
Persistent gender disparities in mathematics continue to influence educational and career choices, often deterring girls from fully engaging with the subject. Prior research, as discussed in \ref{chap:SocialScience} has demonstrated that gender stereotypes and stereotype threat adversely affect both performance and self-efficacy in mathematics. The present study examines whether a brief, school-based intervention can diminish stereotype endorsement and promote greater confidence among female students. Participants, drawn from schools across various educational levels, completed a pre-intervention questionnaire and creative task aimed at assessing both explicit and implicit stereotypes related to mathematics. They then attended a mini-seminar that highlighted the life and contributions of a prominent woman mathematician, followed by a post-intervention questionnaire and a similar activity. Preliminary analyses indicate that exposure to positive female role models in mathematics can reduce the endorsement of gendered beliefs and enhance self-reported comfort with the subject, particularly among girls. These findings underscore the potential of short, accessible interventions to promote gender equity in STEM education and offer practical strategies for educators aiming to counteract the impact of stereotypes in the classroom.

\begin{itemize}
    \item scopo della ricerca e ipotesi: con un piccolo intervento nel quale si parla di una matematica (o simile), quindi portando un modello da seugire femminile, si possono ridurre gli stereotipi/migliorare la self-efficacy delle ragazze. I seguenti fattori influenzano il risultato: età, sesso.
    \item tipologia di intervento: seminario
    \item serve gruppo di controllo?
    \item metodologia
    \item target di partecipanti
    \item questionario: Ragazzi e ragazze sono ugualmente bravi in matematica (1-5) se studio, posso migliorare in matematica (1-5)  Quanto ti senti sicuro in matematica (1-5) sceglierei un corso avanzato se fosse possibile (1-5), Disegna una persona che secondo te è brava in matematica oppure che ha a che fare con i numeri.
    \item discussione risultati
\end{itemize}

